\section{A Conditional Connection to KL-Regularized RL Teachers}
\label{app:rl_teacher_connection}

Our experiments distill RL-trained Qwen3-1.7B teachers into a
Qwen3-1.7B-Base student. A teacher's RL KL reference must be distinguished
from this student initialization. This appendix derives a conditional
connection for ideal KL-regularized RL teachers and a complete-sequence
distillation objective. These assumptions are not established by the
experiments, so the result is a theoretical supplement to the analysis
of domain balance.

\paragraph{Assumptions and sequence objective.}
Let $\mu_d$ denote the fixed KL reference used for teacher $d$, and
let $\pi_{\theta_0}$ denote the student initialization. They need not coincide.
For each domain $d$, fix a prompt distribution $\mathcal D_d$, a
sequence reward $R_d(x,y)$, and a KL penalty coefficient $\beta_d>0$.
Define
\begin{equation}
\begin{aligned}
\mathcal R_d(\theta)
&=\mathbb E_{x\sim\mathcal D_d,\,y\sim\pi_\theta(\cdot\mid x)}
  [R_d(x,y)],\\
K_d(\theta)
&=\mathbb E_{x\sim\mathcal D_d}
  D_{\mathrm{KL}}\!\left(\pi_\theta(\cdot\mid x)
  \|\mu_d(\cdot\mid x)\right),
\qquad J_d(\theta)=\mathcal R_d(\theta)-\beta_d K_d(\theta).
\end{aligned}
\label{eq:rl_teacher_objective}
\end{equation}
Assume compatible support, finite normalizers and expectations, and
regularity sufficient to differentiate these expectations.
Suppose the frozen teacher equals the unrestricted optimum of this
objective over response distributions. The standard KL-RL solution
\citep[Appendix A.1]{rafailov2023dpo} is
\begin{equation}
q_d^*(y\mid x)
=\frac{\mu_d(y\mid x)\exp(R_d(x,y)/\beta_d)}{Z_d(x)},
\qquad
Z_d(x)=\sum_y\mu_d(y\mid x)\exp(R_d(x,y)/\beta_d).
\label{eq:rl_teacher_optimum}
\end{equation}
An RL-trained checkpoint need not attain this optimum. Here $\beta_d$
is the KL penalty coefficient, not a decoding temperature.

\paragraph{Gradient equivalence.}
Consider the complete-sequence reverse KL, without response-length
normalization,
$L_d^{\mathrm{seq}}(\theta)=\mathbb E_{x\sim\mathcal D_d}
D_{\mathrm{KL}}(\pi_\theta(\cdot\mid x)\|q_d^*(\cdot\mid x))$.
Substituting Equation~\ref{eq:rl_teacher_optimum} gives
\begin{equation}
\begin{aligned}
L_d^{\mathrm{seq}}(\theta)
&=K_d(\theta)-\frac{\mathcal R_d(\theta)}{\beta_d}+C_d
=-\frac{J_d(\theta)}{\beta_d}+C_d,\\
-\nabla_\theta L_d^{\mathrm{seq}}(\theta)
&=\frac{1}{\beta_d}\nabla_\theta J_d(\theta),
\qquad C_d=\mathbb E_{x\sim\mathcal D_d}\log Z_d(x).
\end{aligned}
\label{eq:rl_teacher_gradient}
\end{equation}
The constant $C_d$ is independent of the student parameters. Thus the
OPD descent direction equals a scaled RL ascent direction for every
student policy satisfying the assumptions. These are full derivatives,
including the dependence of the response distribution on $\theta$.
This identity does not require the student to be initialized at $\mu_d$.

\paragraph{Implication for domain coefficients.}
For fixed domain weights $\lambda_d$, the first line below gives the
general mixed descent direction. The second line additionally assumes
$\pi_{\theta_0}=\mu_d$ for every domain:
\begin{equation}
\begin{aligned}
-\nabla_\theta\sum_d\lambda_d L_d^{\mathrm{seq}}(\theta)
&=\sum_d\frac{\lambda_d}{\beta_d}\nabla_\theta\mathcal R_d(\theta)
  -\sum_d\lambda_d\nabla_\theta K_d(\theta),\\
\left.-\nabla_\theta\sum_d\lambda_d L_d^{\mathrm{seq}}(\theta)
\right|_{\theta_0}
&=\left.\sum_d\frac{\lambda_d}{\beta_d}
  \nabla_\theta\mathcal R_d(\theta)\right|_{\theta_0}.
\end{aligned}
\label{eq:rl_teacher_mixture}
\end{equation}
The second line uses $\nabla_\theta K_d(\theta_0)=0$ under that additional
assumption. A Base student and a teacher's RL reference cannot be
identified solely by architecture or parameter count; when their
policies differ, the KL-gradient term must be retained even at
initialization. With equal domain weights and comparable reward units,
hypothetical mathematics and coding teachers with $\beta_d=0.05$ and
$0.10$ have reward-gradient coefficients in a $2{:}1$ ratio. This is
neither a gradient-norm ratio nor a capability-gain prediction.
If $R_d=s_d\widetilde R_d$ for $s_d>0$, the coefficient of
$\nabla_\theta\mathbb E[\widetilde R_d]$ is instead
$\lambda_d s_d/\beta_d$. Jointly rescaling reward and $\beta_d$ by the
same positive constant leaves the optimal teacher unchanged.

\paragraph{Relation to the measured token gradients.}
Section~\ref{sec:normalization} analyzes averaging on fixed rollout
batches, independently of how teachers were trained. The present
result shows how teacher construction can introduce an additional
scale in the ideal sequence objective; it does not supply a universal
product of token share, response length, and inverse temperature for
the implemented gradients. Section~\ref{sec:setting} holds sampled
prefixes and lengths fixed when differentiating token losses.
Those local gradients need not equal the complete-sequence derivative
in Equation~\ref{eq:rl_teacher_gradient}; response-dependent length
weights, top-$k$ truncation, and advantage clipping require separate
analysis. No inverse-$\beta_d$ scaling law or $\beta_d$ calibration
benefit is claimed for the teachers or training runs studied here.
